% !TeX root = ../../Book2.tex
% 纲要标题：### 2.2 子模、商模与模同态

\section{子模、商模与模同态}
\label{b2:sec:0202}

模之间的映射若同时保持加法和标量作用，它的核与像便能像在线性代数中那样用来分析映射。要进一步把核压成零，或把像从陪域中模掉，还须知道子模怎样形成商模。这些构造同样适用于非交换环上的左模。

% 纲要要点（供目录核对）：
% * 子模、商模及其泛性质
% * 核、像、余核与同态群；算子交织与自同态环
% * 模的同构定理

\subsection{子模、商模及其泛性质}
\label{b2:subsec:0202:01}

\begin{definition}\label{b2:def:module-homomorphism}
设 \(R\) 是环，\(M,N\) 是左 \(R\) 模。若映射 \(f:M\to N\) 对所有 \(m,m'\in M\) 及 \(r\in R\) 都满足
\[
f(m+m')=f(m)+f(m'),\qquad f(rm)=r\cdot f(m).
\]
就称 \(f\) 为 \(R\) \term[mo-tongtai]{模同态}[module homomorphism]或 \(R\) 线性映射。若模同态是双射，就称它为\term[mo-tonggou]{模同构}[module isomorphism]；存在这种同构时记 \(M\cong N\)。
\end{definition}

同态自动把零元送到零元、把负元送到负元。若 \(f\) 双射，对 \(n=f(m)\)、\(n'=f(m')\)，有 \(f^{-1}(n+n')=m+m'\) 与 \(f^{-1}(rn)=rm\)，故逆映射也是模同态。\(\mathbb Z\) 模同态正是交换群同态，因为可加性迫使它保持全部整数倍。

\begin{proposition}\label{b2:prop:operator-module-intertwiners}
设 \(V,W\) 是域 \(k\) 上的向量空间，\(T\in\operatorname{End}_k(V)\)、\(S\in\operatorname{End}_k(W)\)，相应的算子模为 \(V_T,W_S\)。映射 \(F:V\to W\) 是 \(k[t]\) 模同态，当且仅当它 \(k\) 线性且满足 \(FT=SF\)。因此 \(V_T\cong W_S\) 当且仅当存在满足此等式的 \(k\) 线性同构；在同一个有限维空间上，这等价于两个算子的矩阵相似。
\end{proposition}
\begin{proof}
模同态保持常数多项式的作用，故 \(k\) 线性；保持 \(t\) 的作用恰好给出 \(FT=SF\)。反过来，\cref{b2:prop:polynomial-action-subspace}说明 \(k\) 线性映射满足 \(FT=SF\) 时，对每个 \(p\in k[t]\) 都有 \(Fp(T)=p(S)F\)，所以它保持全部标量作用，是模同态。可逆时等式改写为 \(S=FTF^{-1}\)，得到相似关系。
\end{proof}

由此，第一章的算子相似问题就成为算子模的同构问题；所比较的既是向量空间，也是指定的多项式作用。

\begin{definition}\label{b2:def:submodule}
设 \(R\) 是环，\(M\) 是左 \(R\) 模。若子集 \(L\subseteq M\) 包含 \(0\)，且对任意 \(x,y\in L\)、\(r\in R\) 都有 \(x+y\in L\)、\(-x\in L\) 与 \(rx\in L\)，就称 \(L\) 为 \(M\) 的\term[zimo]{子模}[submodule]。它在限制的运算下本身是左 \(R\) 模，记为 \(L\le M\)。
\end{definition}
\symidx{\(L\le M\)：子模包含记号}

一个非空子集是子模，当且仅当它对 \(x-y\) 与 \(rx\) 封闭：取 \(x=y\) 得零元，取 \(x=0\) 得负元，再由差得到和。左正则模 \({}_RR\) 的子模恰为左理想；不能把这里的左理想自动当作双边理想。对于算子模 \(V_T\)，子模恰为满足 \(T(W)\subseteq W\) 的 \(k\) 线性子空间 \(W\)：一方面它必须对常数与 \(t\) 的作用封闭，另一方面这种封闭性使它对每个 \(T^j\) 和每个 \(f(T)\) 都封闭。

\begin{proposition}\label{b2:prop:quotient-module}
设 \(R\) 是环，\(M\) 是左 \(R\) 模，\(L\le M\)。对 \(m,m'\in M\)、\(r\in R\)，在陪集集合 \(M/L=\{\,m+L\mid m\in M\,\}\) 上定义
\[
(m+L)+(m'+L)=(m+m')+L,\qquad r(m+L)=rm+L.
\]
这给出左 \(R\) 模，称为\term[shangmo]{商模}[quotient module]。商映射 \(q:M\to M/L\)、\(q(m)=m+L\) 是满同态。对任意左 \(R\) 模 \(P\) 及同态 \(g:M\to P\)，存在同态 \(\bar g:M/L\to P\) 满足 \(g=\bar gq\)，当且仅当 \(g(L)=0\)；存在时 \(\bar g\) 唯一。
当 \(g(L)=0\) 时，所要求的分解由交换图表示为
\[
\begin{tikzcd}[column sep=3.5em,row sep=2.2em]
M\arrow[r,"g"]\arrow[d,"q"']&P\\
M/L\arrow[ur,dashed,"\bar g"']&
\end{tikzcd}
\]
\end{proposition}
\symidx{\(M/L\)：模按子模所得的商模}
\begin{proof}
陪集相等 \(m+L=m'+L\) 等价于 \(m-m'\in L\)。如果另取代表元 \(m+\ell\)、\(m'+\ell'\)，两者之和与 \(m+m'\) 的差为 \(\ell+\ell'\in L\)，而 \(r(m+\ell)-rm=r\ell\in L\)。所以两项运算均不依赖代表元。加法的结合律与交换律由 \(M\) 中对应等式下降，零元是 \(L\)，\(m+L\) 的负元是 \(-m+L\)。例如标量乘法的相容性为
\[
(rs)(m+L)=(rs)m+L=r(sm)+L=r\bigl(s(m+L)\bigr),
\]
另外两条分配律和幺性同样由 \(M\) 中的公理得到。公式本身说明 \(q\) 是满同态。

若 \(g=\bar gq\)，则对 \(\ell\in L\)，有 \(g(\ell)=\bar g(0)=0\)。反过来，若 \(g(L)=0\)，定义 \(\bar g(m+L)=g(m)\)。当 \(m-m'\in L\) 时，\(g(m)-g(m')=g(m-m')=0\)，故定义良好。代入商模运算公式，便得到 \(\bar g\) 的可加性与 \(R\) 线性。由于每个陪集都有形式 \(q(m)\)，关系 \(g=\bar gq\) 迫使这个公式成立，也就给出唯一性。
\end{proof}

\subsection{核、像、余核与同态群}
\label{b2:subsec:0202:02}

\begin{definition}\label{b2:def:module-kernel-image-cokernel}
设 \(R\) 是环，\(f:M\to N\) 是左 \(R\) 模之间的同态。定义其\term[he]{核}[kernel]、\term[xiang]{像}[image]及\term[yuhe]{余核}[cokernel]为
\[
\ker f=\{\,m\in M\mid f(m)=0\,\},\qquad
\operatorname{im}f=\{\,f(m)\mid m\in M\,\},
\]
\[
\operatorname{coker}f=N/\operatorname{im}f.
\]
其中核是 \(M\) 的子模，像是 \(N\) 的子模，余核附有商映射 \(q_f:N\to\operatorname{coker}f\)。
\end{definition}
\symidx{\(\operatorname{coker}f\)：模同态 \(f\) 的余核}

这里子模的断言可直接检查：若 \(f(x)=f(y)=0\)，则 \(f(x-y)=0\) 与 \(f(rx)=0\)；若两个元素为 \(f(x),f(y)\)，则其差为 \(f(x-y)\)，\(r\cdot f(x)=f(rx)\)。一个同态是单射，当且仅当核为零，因为 \(f(x)=f(y)\) 等价于 \(x-y\in\ker f\)。它是满射，当且仅当余核为零，因为这恰好表示 \(N=\operatorname{im}f\)。余核把整个像压成零来度量满射性的缺失，并不是从陪域中选出一个与像互补的子模。

\ProseParagraphBreak
记核的包含映射为 \(i_f:\ker f\to M\)。若同态 \(u:P\to M\) 满足 \(fu=0\)，则 \(u(P)\subseteq\ker f\)，故 \(u\) 唯一地分解为 \(u=i_f\bar u\)。余核的商映射满足 \(q_ff=0\)。若同态 \(h:N\to P\) 满足 \(hf=0\)，这正是说 \(h\) 在 \(\operatorname{im}f\) 上为零；由\cref{b2:prop:quotient-module}，它唯一地分解为 \(h=\bar hq_f\)。这两个分解分别由下图的虚线箭头给出：
\[
\begin{tikzcd}[column sep=1.7em,row sep=2em]
P\arrow[r,"u"]\arrow[dr,dashed,"\bar u"']&M\arrow[r,"f"]&N\\
&\ker f\arrow[u,"i_f"']&
\end{tikzcd}
\qquad
\begin{tikzcd}[column sep=1.7em,row sep=2em]
M\arrow[r,"f"]&N\arrow[r,"h"]\arrow[d,"q_f"']&P\\
&\operatorname{coker}f\arrow[ur,dashed,"\bar h"']&
\end{tikzcd}
\]
核与余核在这里都带有指定的映射，分解要求与这些映射相容，而不只是给出一个模的同构类型。

\ProseParagraphBreak
例如同态 \(\mathbb Z\to\mathbb Z\)、\(a\mapsto6a\) 的核为零、像为 \(6\mathbb Z\)、余核为 \(\mathbb Z/6\mathbb Z\)。它是单射，却不是满射。这里 \(6\mathbb Z\) 没有直和补模：任意非零子模 \(L\le\mathbb Z\) 都含有某个非零整数 \(\ell\)，因而 \(6\ell\) 是 \(L\cap6\mathbb Z\) 的非零元素，不能满足直和所要求的交为零。零子模又不能与 \(6\mathbb Z\) 相加得到整个 \(\mathbb Z\)。余核仍然存在，它将全部 \(6\mathbb Z\) 视为零，并不要求从陪域中选出一个补模。

记所有模同态 \(M\to N\) 的集合为 \(\operatorname{Hom}_R(M,N)\)。逐点定义
\[
(f+g)(m)=f(m)+g(m),\qquad (-f)(m)=-f(m).
\]
例如 \((f+g)(rm)=r\cdot f(m)+r\cdot g(m)=r(f+g)(m)\)，所以和仍为模同态；关于加法的其余条件来自 \(N\) 的交换群结构。于是 \(\operatorname{Hom}_R(M,N)\) 总是一个交换群，零元为零映射。
\symidx{\(\operatorname{Hom}_R(M,N)\)：左 \(R\) 模同态组成的交换群}

当定义域与陪域相同，同态还可以复合。记
\[
\operatorname{End}_R(M)\coloneqq\operatorname{Hom}_R(M,M),
\]
以逐点加法和复合乘法 \(fg=f\circ g\) 构成\term[mo-zitongtai-huan]{模自同态环}[endomorphism ring of a module]，单位是恒等映射。复合仍保持加法与 \(R\) 作用，其余环公理与\secref{b2:sec:0201}中的加法自同态环一样逐点核验。

若还要在同态群上定义标量作用，必须检查作用后的映射仍为 \(R\) 线性。新作用可以施加在陪域的值上，也可以先施加在定义域的元素上再取其像；两者都需要与原来的 \(R\) 作用相容。

\begin{proposition}\label{b2:prop:hom-central-action}
设 \(R\) 是环，\(M,N\) 是左 \(R\) 模。环 \(R\) 的中心
\[
Z(R)=\{\,z\in R\mid zr=rz\;\text{对每个}\;r\in R\,\}
\]
通过 \((zf)(m)=z\cdot f(m)\) 使 \(\operatorname{Hom}_R(M,N)\) 成为 \(Z(R)\) 模。特别地，若 \(R\) 交换，同态群以此成为 \(R\) 模。
\end{proposition}
\begin{proof}
中心是含幺交换子环。对 \(z\in Z(R)\)、\(r\in R\)，有
\[
(zf)(rm)=zr\cdot f(m)=rz\cdot f(m)=r\cdot(zf)(m).
\]
因此 \(zf\) 仍为 \(R\) 线性映射；可加性来自 \(N\) 的分配律。对 \(z,z'\in Z(R)\) 逐点计算，\((zz')f=z(z'f)\)、\((z+z')f=zf+z'f\)、\(z(f+g)=zf+zg\)、\(1f=f\)，给出所需模公理。
\end{proof}

不能无条件把 \(Z(R)\) 换成 \(R\)。取 \(M=N={}_RR\)，其中 \(R=M_2(k)\)，令 \(f=\operatorname{id}_R\)，并试图让 \(rf\) 是左乘 \(r\) 的映射。用矩阵单位 \(E_{ij}\) 表示仅第 \((i,j)\) 项为一的矩阵，取 \(r=E_{12}\)、\(s=E_{21}\)，则
\symidx{\(E_{ij}\)：仅指定位置为一的矩阵单位}
\[
(rf)(s)=rs=E_{11},\qquad s(rf)(1_R)=sr=E_{22}.
\]
两者不同，所以 \(rf\) 甚至不是左 \(R\) 模同态。问题在于标量未必能穿过另一个标量，并非加法定义有问题。

\ProseParagraphBreak
设 \(S\) 是环。若陪域 \(N\) 还带左 \(S\) 模结构，且对 \(s\in S\)、\(r\in R\)、\(n\in N\) 满足 \(s(rn)=r(sn)\)，则逐点公式 \((sf)(m)=s\cdot f(m)\) 给 \(\operatorname{Hom}_R(M,N)\) 一个左 \(S\) 模结构。确实，\((sf)(rm)=s(rf(m))=r(sf)(m)\)，所以 \(sf\) 仍为 \(R\) 线性映射；其余模公理由 \(N\) 的作用逐点给出。

\ProseParagraphBreak
另一种左 \(S\) 作用来自定义域。若 \(M\) 同时具有右 \(S\) 模结构，并满足 \((rm)s=r(ms)\)，就称 \(M\) 为 \((R,S)\) \term[shuangmo]{双模}[bimodule]。此时定义
\[
(sf)(m)=f(ms),
\]
这与上一段在 \(N\) 上逐点作用的公式不同。它仍给出左 \(S\) 模结构，因为 \((sf)(rm)=f\bigl((rm)s\bigr)=r\cdot f(ms)\)，而
\[
\bigl(s(tf)\bigr)(m)=(tf)(ms)=f\bigl((ms)t\bigr)=f\bigl(m(st)\bigr)=((st)f)(m).
\]
其余模公理也来自 \(M\) 上的右作用。这里把定义域的右作用变成了同态群上的左作用，次序由 \((ms)t=m(st)\) 保证；它与直接在陪域取值上作用是两种不同的构造。

\begin{proposition}\label{b2:prop:regular-hom-endomorphism}
设 \(R\) 是环，\(M\) 是左 \(R\) 模。取值映射
\[
\operatorname{Hom}_R(R,M)\longrightarrow M,\qquad f\longmapsto f(1_R)
\]
是交换群同构，其逆把 \(m\) 送到 \(r\mapsto rm\)。以复合为乘法的模自同态环满足
\[
\operatorname{End}_R({}_RR)\cong R^{\mathrm{op}}.
\]
\end{proposition}
\begin{proof}
线性条件迫使 \(f(r)=f(r1_R)=r\cdot f(1_R)\)；任意 \(m\) 所给的映射 \(r\mapsto rm\) 确实线性。这两个对应逐点互逆，且保持加法。

当 \(M=R\) 时，\(f\) 因而是某个右乘映射 \(u_a(r)=ra\)。左 \(R\) 线性要求 \(f(sr)=sf(r)\)，参数 \(a=f(1_R)\) 出现在 \(r\) 的右侧；左乘 \(r\mapsto ar\) 则要求 \(asr=sar\)，对一般的 \(a,s\) 并不成立。右乘自同态的复合次序给出
\[
(u_a\circ u_b)(r)=(rb)a=r(ba)=u_{ba}(r).
\]
所以 \(a\mapsto u_a\) 是从反环 \(R^{\mathrm{op}}\) 到该自同态环的保幺环同构。
\end{proof}
\symidx{\(\operatorname{End}_R(M)\)：左 \(R\) 模 \(M\) 的自同态环}

\subsection{模的同构定理}
\label{b2:subsec:0202:03}

\begin{theorem}[第一同构]\label{b2:thm:module-first-isomorphism}
设 \(R\) 是环，\(f:M\to N\) 是左 \(R\) 模之间的同态。它诱导 \(R\) 模同构
\[
M/\ker f\longrightarrow\operatorname{im}f,\qquad m+\ker f\longmapsto f(m).
\]
\end{theorem}
\begin{proof}
把 \(f\) 的陪域改为 \(\operatorname{im}f\)，由\cref{b2:prop:quotient-module}，它在 \(\ker f\) 上为零，因而唯一地下降为所写同态。该同态满射；若 \(f(m)=0\)，则 \(m\in\ker f\)，故源中的陪集为零，所以它也单射。逆映射的线性性由双射同态的一般性质得到。
\end{proof}

这一结果称为\term[mo-diyi-tonggou-dingli]{模的第一同构定理}[first isomorphism theorem for modules]。它所构造的不是 \(M\) 与其像之间任意的同构，而是把全部且仅有的零像元素先模掉。两个向量空间的线性映射满足同一结论，这里不需要维数有限，也不需要标量环是域。

\begin{theorem}[子模对应与商的同构]\label{b2:thm:module-correspondence}
设 \(R\) 是环，\(M\) 是左 \(R\) 模，\(L\le M\)，商映射为 \(q:M\to M/L\)。\(M\) 中包含 \(L\) 的子模 \(U\) 与 \(M/L\) 的子模 \(W\) 通过
\[
U\longmapsto U/L=q(U),\qquad W\longmapsto q^{-1}(W)
\]
一一对应，并保持包含关系。对 \(L\le U\le M\)，有
\[
(M/L)/(U/L)\cong M/U,
\]
同构把 \((m+L)+(U/L)\) 送到 \(m+U\)。
\end{theorem}
\begin{proof}
同态的像与逆像保持子模条件，且 \(q^{-1}(W)\) 总包含 \(L\)。若 \(U\supseteq L\)，则 \(q(m)\in q(U)\) 等价于某个 \(u\in U\) 满足 \(m-u\in L\)，也就等价于 \(m\in U\)。故 \(q^{-1}(q(U))=U\)。由于 \(q\) 满射，对任意 \(W\) 有 \(q(q^{-1}(W))=W\)。两种操作保持包含关系，得到所述对应。

映射 \(M\to M/U\)、\(m\mapsto m+U\) 在 \(L\) 上为零。商映射的泛性质使它下降为满同态 \(M/L\to M/U\)，其核是所有满足 \(m\in U\) 的陪集，即 \(U/L\)。应用\cref{b2:thm:module-first-isomorphism}就得到第二个同构及其元素公式。
\end{proof}

\begin{theorem}[第二同构]\label{b2:thm:module-second-isomorphism}
设 \(R\) 是环，\(M\) 是左 \(R\) 模。若 \(U,V\le M\)，则 \(U+V=\{\,u+v\mid u\in U,\ v\in V\,\}\) 与 \(U\cap V\) 均为子模，并有 \(R\) 模同构
\[
U/(U\cap V)\longrightarrow(U+V)/V,\qquad
u+(U\cap V)\longmapsto u+V.
\]
\end{theorem}
\begin{proof}
和中的两个元素相减后仍为来自 \(U,V\) 的元素之和，标量乘法也保持这种形式；交中的元素同时满足两个子模的封闭条件。因此它们都是子模。考虑同态 \(U\to(U+V)/V\)、\(u\mapsto u+V\)。每个陪集 \((u+v)+V=u+V\)，故它满射；零像的条件为 \(u\in V\)，故核为 \(U\cap V\)。再用\cref{b2:thm:module-first-isomorphism}即可。
\end{proof}

后两个结果分别包含\term[mo-disan-tonggou-dingli]{模的第三同构定理}[third isomorphism theorem for modules]与\term[mo-dier-tonggou-dingli]{模的第二同构定理}[second isomorphism theorem for modules]。两者都先构造一个自然满同态，再计算它的核，最后应用第一同构定理：第二同构使用 \(U\to(U+V)/V\)，第三同构使用 \(M/L\to M/U\)。商模公式记录的是这两个具体映射的结果。作为一个算例，取 \(M=\mathbb Z\)、\(U=4\mathbb Z\)、\(V=6\mathbb Z\)。此时 \(U+V=2\mathbb Z\)、\(U\cap V=12\mathbb Z\)，得到
\[
4\mathbb Z/12\mathbb Z\cong2\mathbb Z/6\mathbb Z.
\]
同构把 \(4+12\mathbb Z\) 送到 \(4+6\mathbb Z\)；两边都有三个元素，但对应公式还告诉我们具体的生成元怎样相配。
